\section{Localization (O2)}
\label{sec:localization}

Detection answers whether something is wrong; localization answers which directed link, which is the question a repair needs. This section is a question about the passive detectors, not about the ledger. Given per-directed-link counts, a detector that flags a link has already localized it. The ledger's row is therefore the oracle against which we measure three things: how far down in $p$ a receive-only vantage can still disalias the uplink and downlink pair, what it costs, and what a counter pair read with skew does to the same information. Figure~\ref{fig:localization} presents the exact-localization rate and the size of the named set for a downlink fault and for an uplink fault, over 50 seeds per cell on the two-hop fabric of Section~\ref{sec:method}; every rate in this section carries a 95 percent Wilson interval of [0.93, 1.00] at 1.00 and [0.00, 0.07] at 0.00, and the full table is in the artifact.

\begin{figure*}[t]
\centering
\includegraphics[width=\textwidth]{fig_localization.pdf}
\caption{Localization under spraying, 50 seeds per cell with 95 percent intervals. (a) Exact-localization rate for a downlink fault. (b) Exact-localization rate for an uplink fault. (c) Mean size of the named set over detections: 1 is the exact directed link, 2 is the uplink and downlink pair; the ledger's set size is 1.00 in every cell, and a missing point means the arm detected nothing.}
\label{fig:localization}
\end{figure*}

\textbf{The oracle row, and the counter pair beside it.} In every cell of both families the ledger's exact rate was 1.00 [0.93, 1.00], its mean set size was exactly 1.00, and it named a healthy link in no seed. This is definitional: it detects in every seed with a false-positive rate of 0.00 (Section~\ref{sec:detection}), and the sublink it rejects is a directed link. The counter pair read at one instant reproduced the row exactly. The counter pair read with the measured best-case skew of $\sigma=0.026$ named a set of about three links in every seed and never the faulty link alone. The reason is that the phantom loss of Section~\ref{sec:detection} makes two or three healthy links look faulty beside the real one; the same information, read on two switches, no longer localizes. The measured result of this section is therefore the passive arms' aliasing and the counter pair's skew, with the ledger as the reference.

\textbf{SprayCheck-Z ties at high loss and degrades to the pair at low loss.} At 1.5 percent in both families and at 1.0 percent for the uplink, SprayCheck's cross-leaf intersection resolved the fault to one directed link as well as the ledger did. There were no discordant seeds between the two arms in those cells, and we report the tie as a null result rather than as a test statistic, since a McNemar test with no discordant pairs is uninformative. At 0.5 percent its exact rate fell to 0.34 for the downlink and 0.26 for the uplink. At $10^{-3}$ it fell to 0.00 and 0.02 with a mean set size of 2.00, i.e., the uplink and downlink pair its paper says an uncorroborated report is left with. The reason is that at those rates its Z-test flags the faulty spine on too few distinct leaves to complete the intersection, so it can name the spine but not the directed link. Every such pair contains one healthy link, which is why the ambiguous column reaches 1.00 there: a repair that acted on the set would touch the wrong hop half the time. At $10^{-4}$ it missed in 74 percent of seeds even under the generous budget.

\textbf{FlowPulse-$\theta$ cannot localize an uplink fault at any rate.} For a downlink fault it was exact in 94 percent of seeds at 1.5 percent, fell to 22 percent at 1.0 percent with a mean set size of 2.74, because its all-or-one rule landed in neither branch and returned the port's whole candidate set, and missed at 0.5 percent and below. For an uplink fault it missed in every seed at every rate, including 1.5 percent. The reason is that a single sender's loss moves the aggregate load of the ingress port by only about $1/\!\text{(number of senders)}$ of the drop rate, here about 0.5 percent for a 1.5 percent uplink loss shared by three senders, which never clears the paper's fixed one-percent threshold. This is faithful to FlowPulse's own design: in its native setting one sender occupies each port and an uplink fault is the whole port's load. Under multi-sender spraying the same fault is diluted below its threshold. The gap is a property of the sprayed fabric, reproduced, not a handicap imposed on the baseline.

\textbf{Boundary.} These results hold for one independent fault on a fabric of four leaves and eight spines with uniform spray. FlowPulse's dilution grows with the number of senders per port, so it is worse at larger radix; SprayCheck's intersection improves with more leaves but needs a longer corroborating round. We expect the ordering of the arms to survive both changes, while the exact crossover rates are properties of this fabric, not universal constants.
